\par\noindent\textbf{(1)} 自由传播距离 \(L\) 的ABCD矩阵为 \(P(L)=\begin{pmatrix}1&L\\0&1\end{pmatrix}\)，曲率半径 \(R\) 的球面镜反射光线矩阵可在一致的方向约定下记为 \(M(R)=\begin{pmatrix}1&0\\-2/R&1\end{pmatrix}\)。往返矩阵由这些基本矩阵按路径顺序相乘，各自行列式均为1。
\par\noindent\textbf{(2)} 周期光线满足线性辛映射，往返矩阵本征值乘积为1。稳定要求本征值位于单位圆，等价于 \(|(A+D)/2|<1\)。对两镜腔算得 \((A+D)/2=2g_1g_2-1\)，故稳定区 \(0<g_1g_2<1\)，端点属于边界退化。
\par\noindent\textbf{(3)} 若一面平镜 \(R_1=\infty\)，另一面凹镜 \(R_2=R>L\)，平镜处波前曲率无限，束腰恰位于平镜。传播到另一镜时 Gaussian曲率 \(R(L)=L[1+(z_R/L)^2]\)。镜面匹配给 \(z_R^2=L(R-L)\)，而 \(w_0^2=\lambda z_R/\pi\)。
\par\noindent\textbf{(4)} 对相同镜面曲率 \(R_1=R_2=L\)，\(g_1g_2=0\) 为共焦稳定边界之一，但光线共振简并显著；而 \(R_j\gg L\) 的近平面腔有 \(g_1g_2\approx1\)，也靠近稳定边界。机械微小失准可能强烈改变光斑位置，严格稳定区应排除边界。
